%%======%%======%0%======%%=======%%
\vspace{5mm}
\section{Uso de IA generativa}
%%======%%======%0%======%%=======%%

O uso de inteligência artificial foi recorrente no processo de desenvolvimento do trabalho. Especial para a solução de bugs, refinamento de códigos e geração de códigos mais extensos ou complexos. A IA foi utilizada em todos os arquivos de códigos presentes no relatório,e teve maior relevância para o desenvolvimento dos códigos em \texttt{javascript} e \texttt{CSS}, uma vez que estes códigos apresentam maior rigidez, maior complexidade e macetes que são indispensáveis para boas métricas da interface desenvolvida. 

Outro uso relevante de IA no trabalho, foi voltado para a constante validação dos avanços realizados, sendo que a ferramenta, auxiliou o desenvolvimento do projeto fazendo comparações entre a estrutura em desenvolvimento e as exigências do professor referente ao trabalho 1. Desse modo, a IA teve acesso ao código constantemente. 

No que tange ao HTML, a IA trouxe caracteristicas marcantes com o uso de ARIA, ferramentas como o GPT e o CLAUDE inserem com alta frequência o recurso, sendo em alguns momentos, necessário o controle e a reestruturação, prevenindo redundâncias ou excessos por parte do código.

Deve ser ressaltado, que a inteligência artificial (especificamente a ferramenta Claude), desenvolveu um modelo de site, baseado nos preceitos discutidos anteriormente, o qual foi usado como esqueleto principal para o desenvolvimento da interface atual. Porém, elementos como \texttt{javascript} e \texttt{CSS} sofreram alterações quase que totais, uma vez que, ao longo do processo, novas implementações, readequações e certos problemas ou soluções surgiram.

Os códigos, todos, passaram por revisão e refinamento de inteligência artificial, entretanto, os trechos que foram gerados significativamente por IA, estão claramente indicados no código. Destacam-se os arquivos \texttt{components.css}, devido a sua extensão e abrangência, \texttt{i18n.js} devido a sua complexidade, a IA foi especialmente usada para criar a lista de palavras a serem traduzidas e \texttt{adicionar-trabalho.js}, pois foi encontrada uma dificuldade relevante pelo aluno autor de fazer o código funcionar devidamente, vale destacar que o \texttt{README} foi enfeitado pelo Gemini PRO, entretanto o texto ali presente é INTEGRALMENTE escrito pelo aluno autor. Foi apenas corrigido e estruturado mais estéticamente em \texttt{markdown}.

Dessa forma, a inteligência artificial foi utilizada como ferramenta de apoio durante o desenvolvimento, revisão, estruturação e validação do projeto, enquanto as decisões sobre a implementação, alterações no código e adequação às necessidades do trabalho permaneceram sob controle do desenvolvedor. 

Uma tabela (Tabela \ref{tab:uso-ia-arquivos}) mais detalhada pode ser consultado na Seção \ref{sec:apendice}, de apêndices.
